\documentclass[../main.tex]{subfiles}

\begin{document}


\section{Kiến trúc tổng thể}

\subsection{Tổng quan kiến trúc}

Hệ thống PPAT được xây dựng theo kiến trúc hướng dịch vụ (Service-Oriented Architecture - SOA), trong đó mỗi thành phần đảm nhận một trách nhiệm riêng biệt và giao tiếp với nhau qua các giao thức được định nghĩa rõ ràng. Cách tổ chức này giúp các thành phần có thể phát triển, nâng cấp và mở rộng độc lập mà không ảnh hưởng đến phần còn lại của hệ thống.

\begin{figure}[H]
\centering
\includegraphics[width=1.0\textwidth]{bieu_do_thanh_phan_kien_truc_soa.png}
\caption{Biểu đồ thành phần kiến trúc SOA}
\label{fig:bieu-do-thanh-phan-kien-truc-soa}
\end{figure}

Như thể hiện trong Hình 3.1, toàn bộ hệ thống gồm năm thành phần chính: Frontend, Backend API, LLM Service, Code Runner Service và tầng dữ liệu (PostgreSQL + Redis).

Frontend được xây dựng bằng Next.js tích hợp Monaco Editor, là điểm tiếp xúc trực tiếp với người dùng. Trình duyệt giao tiếp với Frontend qua HTTP/HTTPS và nhận về HTML/CSS/JS để render giao diện. Frontend không xử lý logic nghiệp vụ mà chỉ đảm nhận việc hiển thị và thu thập tương tác từ sinh viên, bao gồm giao diện học bài, code editor, chat với AI Tutor và trang theo dõi tiến độ.

Backend API sử dụng Fastify đóng vai trò trung tâm điều phối toàn bộ hệ thống. Đây là nơi xử lý xác thực JWT, quản lý người dùng, bài học, bài tập và hội thoại. Mọi request từ Frontend đều đi qua Backend qua REST API. Backend cũng là nơi duy nhất kết nối trực tiếp với cả PostgreSQL, Redis, LLM Service và Code Runner Service đảm bảo logic nghiệp vụ được tập trung, không phân tán.

LLM Service hoạt động như một AI Orchestrator độc lập. Khi nhận yêu cầu từ Backend, service này truy vấn knowledge log của sinh viên, xây dựng dynamic system prompt phù hợp với từng cá nhân, sau đó gọi OpenAI API (GPT-4o) qua HTTPS. Phản hồi từ GPT-4o được stream ngược về Backend rồi tiếp tục forward về Frontend thông qua Server-Sent Events (SSE) đây là giao tiếp không đồng bộ duy nhất trong hệ thống, được biểu diễn bằng mũi tên nét đứt màu xanh trong Hình 3.1.

Code Runner Service là thành phần xử lý việc thực thi code của sinh viên. Khi sinh viên nhấn chạy code, Backend gửi yêu cầu đến service này qua REST API. Code được chạy bên trong một Docker Sandbox hoàn toàn cô lập không có quyền truy cập mạng, giới hạn tài nguyên CPU/RAM và timeout cứng nhằm đảm bảo an toàn cho server. Kết quả thực thi (output hoặc error) được trả về Backend rồi hiển thị cho sinh viên.

Tầng dữ liệu gồm PostgreSQL đảm nhận lưu trữ bền vững toàn bộ dữ liệu nghiệp vụ (người dùng, bài học, hội thoại, tiến độ, analytics) và Redis đảm nhận cache, session và rate limiting. Backend giao tiếp với PostgreSQL qua SQL và với Redis qua Redis Protocol (TCP).

\subsection{Luồng xử lý theo loại request}

Tùy vào loại tương tác của sinh viên, hệ thống xử lý theo các luồng khác nhau:


{
\renewcommand{\arraystretch}{1.5}
\begin{tabularx}{\textwidth}{|p{3.5cm}|X|}
\caption{Ví dụ luồng xử lý theo loại request}
\label{tab:vi-du-luong-xu-ly-theo-loai-request}
\\
\hline
\textbf{Loại request} & \textbf{Luồng xử lý} \\
\hline
Đăng nhập & Frontend → Backend → PostgreSQL (xác thực) → Redis (lưu session) → trả JWT \\
\hline
Xem bài học & Frontend → Backend → Redis (cache check) → PostgreSQL (nếu cache miss) → trả dữ liệu \\
\hline
Chat AI & Frontend → Backend → LLM Service → OpenAI API → stream SSE ngược về Frontend \\
\hline
Chạy code & Frontend → Backend → Code Runner → Docker Sandbox → trả output/error \\
\hline
Xem tiến độ & Frontend → Backend → Redis (cache) → PostgreSQL (aggregate query) → trả dashboard \\
\hline
\end{tabularx}
}

Thiết kế theo luồng này đảm bảo các request thông thường (xem bài học, xem tiến độ) được phục vụ nhanh nhờ Redis cache, trong khi các request nặng hơn (chat AI, chạy code) được tách ra xử lý ở service riêng mà không ảnh hưởng đến hiệu năng tổng thể của Backend.

\section{Thiết kế API}

Hệ thống cung cấp 14 endpoint REST API, được nhóm theo chức năng. Toàn bộ các endpoint yêu cầu xác thực đều kiểm tra JWT token trong header Authorization: Bearer \&lt;token\&gt;.


{
\renewcommand{\arraystretch}{1.5}
\begin{tabularx}{\textwidth}{|p{2.0cm}|X|p{2.0cm}|X|}
\caption{Danh sách endpoint REST API}
\label{tab:danh-sach-endpoint-rest-api}
\\
\hline
\textbf{Method} & \textbf{Endpoint} & \textbf{Phân quyền} & \textbf{Mô tả} \\
\hline
POST & /\allowbreak api/\allowbreak auth/\allowbreak register & Public & Đăng ký tài khoản mới \\
\hline
POST & /\allowbreak api/\allowbreak auth/\allowbreak login & Public & Đăng nhập, trả về JWT và refresh token \\
\hline
GET & /\allowbreak api/\allowbreak lessons & Student & Danh sách bài học, hỗ trợ filter theo ngôn ngữ, chủ đề, độ khó \\
\hline
GET & /\allowbreak api/\allowbreak lessons/\allowbreak :id & Student & Chi tiết một bài học \\
\hline
GET & /\allowbreak api/\allowbreak exercises/\allowbreak :id & Student & Chi tiết bài tập kèm starter code \\
\hline
POST & /\allowbreak api/\allowbreak exercises/\allowbreak :id/\allowbreak submit & Student & Nộp code, chạy test cases, trả kết quả pass/fail/error \\
\hline
GET & /\allowbreak api/\allowbreak conversations & Student & Danh sách phiên hội thoại của sinh viên \\
\hline
POST & /\allowbreak api/\allowbreak conversations & Student & Tạo phiên hội thoại mới \\
\hline
POST & /\allowbreak api/\allowbreak conversations/\allowbreak :id/\allowbreak messages & Student & Gửi tin nhắn, nhận phản hồi AI qua SSE stream \\
\hline
GET & /\allowbreak api/\allowbreak progress/\allowbreak me & Student & Dashboard tiến độ cá nhân \\
\hline
GET & /\allowbreak api/\allowbreak analytics/\allowbreak me & Student & Knowledge log, điểm mạnh và điểm yếu \\
\hline
GET & /\allowbreak api/\allowbreak admin/\allowbreak lessons & Admin & Danh sách bài học phía quản trị \\
\hline
POST & /\allowbreak api/\allowbreak admin/\allowbreak lessons & Admin & Tạo bài học mới \\
\hline
GET & /\allowbreak api/\allowbreak admin/\allowbreak reports & Admin & Báo cáo tổng hợp toàn bộ sinh viên \\
\hline
\end{tabularx}
}

Endpoint quan trọng nhất về mặt kỹ thuật là POST /\allowbreak api/\allowbreak conversations/\allowbreak :id/\allowbreak messages. Khác với các endpoint thông thường trả về JSON, endpoint này trả về response theo định dạng text/event-stream, tức là giữ kết nối mở và đẩy từng token từ GPT-4o về client theo thời gian thực thay vì chờ sinh xong toàn bộ câu trả lời. Cơ chế streaming này được mô tả chi tiết hơn ở mục 3.4.4.

\section{Thiết kế cơ chế Cá nhân hóa AI}

\subsection{Pipeline xử lý tin nhắn (Message Lifecycle)}

Đây là phần kỹ thuật cốt lõi tạo ra sự khác biệt so với việc dùng ChatGPT thông thường. Khi sinh viên gửi một tin nhắn, hệ thống không chuyển thẳng nội dung đó đến GPT-4o mà thực hiện một pipeline 5 bước như mô tả trong Hình 3.3.

Bước 1 - Sinh viên gửi tin nhắn: Frontend gửi POST request đến /\allowbreak api/\allowbreak conversations/\allowbreak :id/\allowbreak messages và đồng thời mở kết nối EventSource để lắng nghe SSE stream. Ở bước này, ba sinh viên với ba trình độ khác nhau gửi cùng một câu hỏi sự khác biệt trong phản hồi bắt đầu từ bước tiếp theo.

Bước 2 - Fetch Student Profile: Backend truy vấn PostgreSQL để lấy toàn bộ thông tin học tập của sinh viên, bao gồm trình độ hiện tại, danh sách chủ đề đã nắm, các lỗi thường gặp từ user\_analytics và lịch sử 10 tin nhắn gần nhất. Đây là bước hoàn toàn không có trong ChatGPT thông thường mỗi sinh viên trả về một tập dữ liệu khác nhau dù cùng gửi một câu hỏi.

Bước 3 - Build Dynamic System Prompt: AI Orchestrator nhận dữ liệu từ bước 2 và điền vào template system prompt có sẵn. Kết quả là ba sinh viên khác nhau tạo ra ba system prompt hoàn toàn khác nhau trước khi gửi đến GPT-4o. Ví dụ với sinh viên trình độ intermediate hay quên base case, system prompt sẽ hướng AI dùng chiến lược bridge từ vòng lặp sang đệ quy và chủ động nhắc về điều kiện dừng.

Bước 4 - Gọi GPT-4o với stream: Backend gọi OpenAI SDK với tham số stream: true. Từng token được GPT-4o sinh ra được forward ngay về Frontend qua Server-Sent Events thay vì chờ đến khi hoàn thành toàn bộ câu trả lời. Sinh viên thấy phản hồi xuất hiện dần dần trên màn hình, tạo cảm giác tương tác tự nhiên như đang chat với người thật.

Bước 5 - Lưu message và cập nhật log: Sau khi stream kết thúc, một background job được kích hoạt để INSERT tin nhắn vào bảng messages và UPDATE user\_analytics cập nhật confidence score, ghi nhận lỗi mới nếu có, và tính toán gợi ý bài học tiếp theo. Bước này chạy bất đồng bộ nên không làm chậm trải nghiệm của sinh viên.

\begin{figure}[H]
\centering
\includegraphics[width=1.0\textwidth]{pipeline_xu_ly_tin_nhan_message_lifecycle.png}
\caption{Pipeline xử lý tin nhắn (Message Lifecycle)}
\label{fig:pipeline-xu-ly-tin-nhan-message-lifecycle}
\end{figure}

\subsection{Dynamic System Prompt}

Cơ chế dynamic system prompt là yếu tố kỹ thuật trực tiếp tạo ra sự cá nhân hóa trong phản hồi của AI. Thay vì dùng một system prompt cố định cho mọi người dùng, hệ thống xây dựng prompt riêng cho từng sinh viên trước mỗi lần gọi GPT-4o.

Template prompt được thiết kế gồm hai phần. Phần đầu là profile data các placeholder được điền động từ database:

\begin{itemize}
    \item \textit{You are a personalized programming tutor AI.}
    \item \textit{Here is the student's profile:}
    \item \textit{Name: {student\_name}}
    \item \textit{Current level: {level} (beginner / intermediate / advanced)}
    \item \textit{Languages learning: {languages}}
    \item \textit{Topics already mastered: {known\_topics}}
    \item \textit{Common mistakes: {error\_patterns}}
    \item \textit{Lessons completed: {completed\_lessons}/{total\_lessons}}
    \item \textit{Recent interaction summary: {recent\_summary}}
\end{itemize}

Các placeholder này được lấy từ các bảng khác nhau trong database: {student\_name} và {level} từ bảng users, {known\_topics} và {error\_patterns} từ cột JSONB trong userAnalytics, {completed\_lessons} tổng hợp từ userProgress, còn {recent\_summary} là tóm tắt từ 10 tin nhắn gần nhất trong messages.

Phần thứ hai là teaching guidelines bộ quy tắc cố định áp dụng cho mọi sinh viên, định nghĩa hành vi của AI Tutor:

\begin{itemize}
    \item \textit{Teaching guidelines:}
    \item \textit{1. NEVER give the full solution directly. Use the Socratic method - ask guiding questions.}
    \item \textit{2. Adjust explanation depth to the student's level.}
    \item \textit{3. If the student hasn't learned a prerequisite, gently redirect to that first.}
    \item \textit{4. For intermediate/advanced students, suggest edge cases, best practices, or comparisons.}
    \item \textit{5. Keep responses concise. Use code blocks for examples.}
    \item \textit{6. Acknowledge when the student does something right - build confidence.}
    \item \textit{7. If the student is stuck after 2 attempts, provide a small hint, not the answer.}
\end{itemize}

Bảy quy tắc này phản ánh trực tiếp các nguyên tắc sư phạm đã trình bày ở Chương 2: không đưa đáp án trực tiếp (quy tắc 1), điều chỉnh theo trình độ (quy tắc 2), phát hiện kiến thức tiên quyết còn thiếu (quy tắc 3), và chỉ gợi ý nhỏ sau hai lần thất bại (quy tắc 7). Nhờ phần profile data thay đổi theo từng sinh viên, cùng một bộ quy tắc này cho ra hành vi hoàn toàn khác nhau sinh viên beginner được giải thích bằng analogy đơn giản, sinh viên advanced được thách thức thêm với edge case và so sánh độ phức tạp thuật toán.

\subsection{Adaptive Learning Path}

Hệ thống không chỉ cá nhân hóa ở phần chat mà còn tự động điều chỉnh lộ trình học tập dựa trên kết quả làm bài tập của sinh viên. Mỗi lần nộp bài, thay vì chỉ trả về pass/fail đơn thuần, hệ thống phân nhánh xử lý theo 5 trường hợp.

\begin{figure}[H]
\centering
\includegraphics[width=1.0\textwidth]{adaptive_learning_path.png}
\caption{Adaptive Learning Path}
\label{fig:adaptive-learning-path}
\end{figure}

Khi sinh viên pass toàn bộ test case, hệ thống ghi nhận milestone vào userProgress, cập nhật knownTopics trong userAnalytics và mở khóa bài tập tiếp theo. AI Tutor sẽ gửi lời chúc mừng ngắn và gợi ý bài học tiếp theo phù hợp với lộ trình.

Khi fail lần đầu, hệ thống ghi nhận loại lỗi vào errorPatterns nhưng không hiển thị gì thêm trên giao diện. Thông tin này được dùng ở phiên chat tiếp theo AI sẽ chủ động nhắc về lỗi đó mà sinh viên không cần phải hỏi lại.

Khi fail lần hai, một gợi ý nhỏ xuất hiện ngay trong giao diện code editor. Gợi ý này chỉ hướng đến điểm sai cụ thể chứ không đưa ra đáp án ví dụ "kiểm tra lại điều kiện dừng của vòng lặp" thay vì viết sẵn code đúng.

Khi fail từ lần ba trở đi, hệ thống tra cứu prerequisiteLessonId của bài tập hiện tại để xác định kiến thức tiên quyết còn thiếu. Thay vì để sinh viên tiếp tục loay hoay, AI Tutor đề xuất ôn lại bài học nền tảng trước và giải thích nhẹ nhàng lý do cần làm vậy.

Khi xảy ra runtime error, hệ thống lưu lại stack trace và chuyển thông tin này vào context của AI Tutor. Thay vì đưa ra fix trực tiếp, AI hướng dẫn sinh viên cách đọc thông báo lỗi và tự xác định nguyên nhân đây là kỹ năng debug thực tế quan trọng hơn việc biết đáp án của một bài cụ thể.

Toàn bộ các nhánh xử lý trên đều hội tụ về một điểm chung là cập nhật userAnalytics, đảm bảo knowledge log luôn phản ánh trạng thái học tập mới nhất của sinh viên và system prompt ở lần chat tiếp theo sẽ chính xác hơn.

\subsection{Streaming Architecture (SSE)}

Thay vì chờ GPT-4o sinh xong toàn bộ câu trả lời rồi mới trả về, hệ thống sử dụng Server-Sent Events (SSE) để forward từng token về phía client ngay khi có. Cách tiếp cận này giúp sinh viên thấy phản hồi xuất hiện dần dần theo thời gian thực, tạo trải nghiệm tự nhiên hơn so với chờ đợi màn hình trắng.

Luồng xử lý diễn ra theo thứ tự sau. Khi sinh viên gửi tin nhắn, Frontend gọi POST request đến endpoint chat và ngay lập tức mở một EventSource connection để lắng nghe stream. Phía Backend, sau khi hoàn thành bước build dynamic system prompt, gọi OpenAI SDK với stream: true và set header response là Content-Type: text/event-stream để giữ kết nối mở. Mỗi token được GPT-4o trả về, Backend forward ngay qua SSE mà không buffer lại. Frontend lắng nghe từng SSE event, lấy nội dung token và append vào chat bubble hiện tại theo từng ký tự.

Khi GPT-4o gửi tín hiệu kết thúc \[DONE\], Backend đóng SSE connection và kích hoạt bất đồng bộ hai tác vụ: INSERT toàn bộ nội dung tin nhắn vào bảng messages và UPDATE userAnalytics. Việc tách phần lưu dữ liệu ra khỏi luồng streaming là có chủ ý sinh viên nhận được token đầu tiên trong vòng vài trăm millisecond mà không bị block bởi các tác vụ ghi database phía sau.

Một điểm cần lưu ý là SSE là giao thức một chiều từ server về client, khác với WebSocket hai chiều. Điều này phù hợp với bài toán chat AI vì trong khi AI đang trả lời, sinh viên không cần gửi thêm dữ liệu chỉ cần nhận stream. Khi sinh viên muốn gửi tin nhắn mới, một POST request hoàn toàn mới được khởi tạo và một SSE stream mới được mở.

\section{Thiết kế chi tiết}

\subsection{Đăng ký}

\begin{figure}[H]
\centering
\includegraphics[width=1.0\textwidth]{giao_dien_chuc_nang_dang_ky.png}
\caption{Giao diện chức năng đăng ký}
\label{fig:giao-dien-chuc-nang-dang-ky}
\end{figure}

\begin{figure}[H]
\centering
\includegraphics[width=1.0\textwidth]{so_do_hoat_dong_chuc_nang_dang_ky.png}
\caption{Sơ đồ hoạt động chức năng Đăng ký}
\label{fig:so-do-hoat-dong-chuc-nang-dang-ky}
\end{figure}

\begin{figure}[H]
\centering
\includegraphics[width=1.0\textwidth]{so_do_trinh_tu_chuc_nang_dang_ky.png}
\caption{Sơ đồ trình tự chức năng Đăng ký}
\label{fig:so-do-trinh-tu-chuc-nang-dang-ky}
\end{figure}

\subsection{Đăng nhập}

\begin{figure}[H]
\centering
\includegraphics[width=1.0\textwidth]{giao_dien_chuc_nang_dang_nhap.png}
\caption{Giao diện chức năng Đăng nhập}
\label{fig:giao-dien-chuc-nang-dang-nhap}
\end{figure}

\begin{figure}[H]
\centering
\includegraphics[width=1.0\textwidth]{so_do_hoat_dong_chuc_nang_dang_nhap.png}
\caption{Sơ đồ hoạt động chức năng Đăng nhập}
\label{fig:so-do-hoat-dong-chuc-nang-dang-nhap}
\end{figure}

\begin{figure}[H]
\centering
\includegraphics[width=1.0\textwidth]{so_do_trinh_tu_chuc_nang_dang_nhap.png}
\caption{Sơ đồ trình tự chức năng Đăng nhập}
\label{fig:so-do-trinh-tu-chuc-nang-dang-nhap}
\end{figure}

\subsection{Dashboard kết hợp Theo dõi tiến độ}

\begin{figure}[H]
\centering
\includegraphics[width=1.0\textwidth]{giao_dien_dashboard_ket_hop_theo_doi_tien_do.png}
\caption{Giao diện Dashboard kết hợp theo dõi tiến độ}
\label{fig:giao-dien-dashboard-ket-hop-theo-doi-tien-do}
\end{figure}

\begin{figure}[H]
\centering
\includegraphics[width=1.0\textwidth]{so_do_hoat_dong_chuc_nang_xem_dashboard_ket_hop_theo_doi_tien_do.png}
\caption{Sơ đồ hoạt động chức năng xem Dashboard kết hợp theo dõi tiến độ}
\label{fig:so-do-hoat-dong-chuc-nang-xem-dashboard-ket-hop-theo-doi-tien-do}
\end{figure}

\begin{figure}[H]
\centering
\includegraphics[width=1.0\textwidth]{so_do_trinh_tu_chuc_nang_xem_dashboard_ket_hop_theo_doi_tien_do.png}
\caption{Sơ đồ trình tự chức năng xem Dashboard kết hợp theo dõi tiến độ}
\label{fig:so-do-trinh-tu-chuc-nang-xem-dashboard-ket-hop-theo-doi-tien-do}
\end{figure}

\subsection{Danh sách Bài học (Xem và Lọc bài học)}

\begin{figure}[H]
\centering
\includegraphics[width=1.0\textwidth]{giao_dien_danh_sach_bai_hoc.png}
\caption{Giao diện Danh sách bài học}
\label{fig:giao-dien-danh-sach-bai-hoc}
\end{figure}

\begin{figure}[H]
\centering
\includegraphics[width=1.0\textwidth]{so_do_hoat_dong_chuc_nang_xem_danh_sach_bai_hoc.png}
\caption{Sơ đồ hoạt động chức năng xem Danh sách bài học}
\label{fig:so-do-hoat-dong-chuc-nang-xem-danh-sach-bai-hoc}
\end{figure}

\begin{figure}[H]
\centering
\includegraphics[width=1.0\textwidth]{so_do_trinh_tu_chuc_nang_xem_danh_sach_bai_hoc.png}
\caption{Sơ đồ trình tự chức năng xem Danh sách bài học}
\label{fig:so-do-trinh-tu-chuc-nang-xem-danh-sach-bai-hoc}
\end{figure}

\subsection{Chi tiết Bài học}

\begin{figure}[H]
\centering
\includegraphics[width=1.0\textwidth]{giao_dien_chi_tiet_bai_hoc.png}
\caption{Giao diện Chi tiết bài học}
\label{fig:giao-dien-chi-tiet-bai-hoc}
\end{figure}

\begin{figure}[H]
\centering
\includegraphics[width=1.0\textwidth]{so_do_hoat_dong_chuc_nang_xem_chi_tiet_bai_hoc.png}
\caption{Sơ đồ hoạt động chức năng xem Chi tiết bài học}
\label{fig:so-do-hoat-dong-chuc-nang-xem-chi-tiet-bai-hoc}
\end{figure}

\begin{figure}[H]
\centering
\includegraphics[width=1.0\textwidth]{so_do_trinh_tu_chuc_nang_xem_chi_tiet_bai_hoc.png}
\caption{Sơ đồ trình tự chức năng xem Chi tiết bài học}
\label{fig:so-do-trinh-tu-chuc-nang-xem-chi-tiet-bai-hoc}
\end{figure}

\subsection{Code Editor (Làm bài tập)}

\begin{figure}[H]
\centering
\includegraphics[width=1.0\textwidth]{giao_dien_code_editor_lam_bai_tap.png}
\caption{Giao diện Code Editor (Làm bài tập)}
\label{fig:giao-dien-code-editor-lam-bai-tap}
\end{figure}

\begin{figure}[H]
\centering
\includegraphics[width=1.0\textwidth]{so_do_hoat_dong_chuc_nang_code_editor_lam_bai_tap.png}
\caption{Sơ đồ hoạt động chức năng Code editor (Làm bài tập)}
\label{fig:so-do-hoat-dong-chuc-nang-code-editor-lam-bai-tap}
\end{figure}

\begin{figure}[H]
\centering
\includegraphics[width=1.0\textwidth]{so_do_trinh_tu_chuc_nang_code_editor_lam_bai_tap.png}
\caption{Sơ đồ trình tự chức năng Code editor (Làm bài tập)}
\label{fig:so-do-trinh-tu-chuc-nang-code-editor-lam-bai-tap}
\end{figure}

\subsection{Chat với AI Tutor}

\begin{figure}[H]
\centering
\includegraphics[width=1.0\textwidth]{giao_dien_chat_voi_ai_tutor.png}
\caption{Giao diện Chat với AI Tutor}
\label{fig:giao-dien-chat-voi-ai-tutor}
\end{figure}

\begin{figure}[H]
\centering
\includegraphics[width=1.0\textwidth]{giao_dien_chat_ai_voi_tutor_trong_trinh_editor.png}
\caption{Giao diện Chat AI với Tutor trong trình Editor}
\label{fig:giao-dien-chat-ai-voi-tutor-trong-trinh-editor}
\end{figure}

\begin{figure}[H]
\centering
\includegraphics[width=1.0\textwidth]{so_do_hoat_dong_chuc_nang_chat_ai_voi_tutor.png}
\caption{Sơ đồ hoạt động chức năng Chat AI với Tutor}
\label{fig:so-do-hoat-dong-chuc-nang-chat-ai-voi-tutor}
\end{figure}

\begin{figure}[H]
\centering
\includegraphics[width=1.0\textwidth]{so_do_trinh_tu_chuc_nang_chat_ai_voi_tutor.png}
\caption{Sơ đồ trình tự chức năng Chat AI với Tutor}
\label{fig:so-do-trinh-tu-chuc-nang-chat-ai-voi-tutor}
\end{figure}

\section{Thiết kế Cơ sở dữ liệu}

\subsection{Sơ đồ lớp thực thể (ERD)}

\begin{figure}[H]
\centering
\includegraphics[width=1.0\textwidth]{so_do_lop_thuc_the_he_thong.png}
\caption{Sơ đồ lớp thực thể hệ thống}
\label{fig:so-do-lop-thuc-the-he-thong}
\end{figure}

Cơ sở dữ liệu của hệ thống PPAT được thiết kế trên PostgreSQL gồm 8 bảng chính, thể hiện trong Hình 3.26. Các bảng được tổ chức xoay quanh thực thể trung tâm là Users mọi dữ liệu học tập, hội thoại và tiến độ đều gắn với một người dùng cụ thể.

Quan hệ giữa các bảng được xác định như sau: một người dùng có thể có nhiều cuộc hội thoại (1:n giữa Users và Conversations), mỗi hội thoại chứa nhiều tin nhắn (1:n giữa Conversations và Messages). Về phía học tập, mỗi bài học có thể có nhiều bài tập (1:n giữa Lessons và Exercises), và mỗi lần sinh viên nộp bài được ghi nhận riêng trong ExerciseSubmissions. Tiến độ học tập của từng sinh viên trên từng bài học và bài tập được theo dõi qua UserProgress, trong khi UserAnalytics lưu trữ knowledge log bao gồm các chủ đề đã nắm, lỗi thường gặp và ước tính trình độ hiện tại phục vụ cho cơ chế cá nhân hóa của AI Tutor.

\subsection{Mô tả chi tiết các bảng}

Bảng Users là trung tâm của toàn bộ hệ thống, lưu thông tin tài khoản bao gồm email, mật khẩu đã được hash, tên hiển thị, và hai trường enum quan trọng là level (trình độ: beginner, intermediate, advanced) và role (phân quyền: student, admin). Trường level được cập nhật tự động khi sinh viên hoàn thành đủ số bài tập theo ngưỡng quy định.

Bảng Conversations và Messages lưu toàn bộ lịch sử hội thoại giữa sinh viên và AI Tutor. Mỗi Message ghi nhận vai trò người gửi qua trường role (kiểu enum messageRole: user hoặc assistant), nội dung tin nhắn và số token đã tiêu thụ thông tin hữu ích để theo dõi chi phí API.

Bảng Lessons lưu nội dung bài học dưới dạng Markdown trong trường contentMd, phân loại theo ngôn ngữ lập trình (language), chủ đề (topic) và độ khó (difficulty). Thiết kế này cho phép lọc và gợi ý bài học phù hợp với trình độ từng sinh viên một cách linh hoạt.

Bảng Exercises lưu thông tin bài tập gắn với từng bài học qua lessonId. Điểm đáng chú ý là trường testCases có kiểu JSONB, cho phép lưu linh hoạt nhiều bộ test case với cấu trúc input/output khác nhau mà không cần thêm bảng riêng. Trường prerequisiteLessonId cho phép khai báo bài học điều kiện nếu sinh viên chưa hoàn thành bài học tiên quyết, AI Tutor sẽ chủ động chuyển hướng thay vì giải thích ngay.

Bảng ExerciseSubmissions ghi nhận mỗi lần nộp bài của sinh viên, bao gồm code đã nộp, trạng thái kết quả (status kiểu enum submissionStatus: pass, fail, error), thông báo lỗi nếu có và thời gian thực thi tính bằng millisecond. Dữ liệu này là đầu vào chính cho cơ chế phát hiện lỗi thường gặp trong knowledge log.

Bảng UserProgress theo dõi trạng thái hoàn thành của từng sinh viên trên từng bài học và bài tập, lưu điểm số và trạng thái hoàn thành. Bảng này được truy vấn tổng hợp khi render dashboard tiến độ.

Bảng UserAnalytics là bảng quan trọng nhất đối với tính năng cá nhân hóa. Hai trường JSONB knownTopics và errorPatterns lưu trữ dạng key-value linh hoạt ví dụ errorPatterns có thể chứa {"missing\_base\_case": 3, "off\_by\_one": 2} để AI biết sinh viên hay mắc lỗi gì và mức độ thường xuyên. Trường levelEstimate là trình độ ước tính theo thời gian thực, có thể khác với level chính thức trong bảng users và được dùng để điều chỉnh độ khó phản hồi của AI.


{
\renewcommand{\arraystretch}{1.5}
\begin{tabularx}{\textwidth}{|p{4cm}|X|p{4cm}|}
\caption{Các bảng thực thể trong hệ thống}
\label{tab:cac-bang-thuc-the-trong-he-thong}
\\
\hline
\textbf{Bảng} & \textbf{Các cột chính} & \textbf{Mục đích} \\
\hline
users & id, email, password\_hash, display\_name, level, created\_at & Tài khoản người dùng \\
\hline
lessons & id, title, language, topic, difficulty, content\_md, order\_index, created\_by & Bài học lý thuyết \\
\hline
exercises & id, lesson\_id, title, description, starter\_code, expected\_output, test\_cases (JSONB), difficulty & Bài tập thực hành \\
\hline
exercise\_submissions & id, user\_id, exercise\_id, code, status, error\_msg, runtime\_ms, submitted\_at & Lịch sử nộp bài \\
\hline
user\_progress & id, user\_id, lesson\_id, exercise\_id, completion\_status, score, completed\_at & Tiến độ từng bài \\
\hline
user\_analytics & id, user\_id, known\_topics (JSONB), error\_patterns (JSONB), level\_estimate, last\_updated & Knowledge log cá nhân hóa \\
\hline
conversations & id, user\_id, title, created\_at, last\_message\_at & Phiên hội thoại với AI \\
\hline
messages & id, conversation\_id, role (user/assistant), content, tokens\_used, created\_at & Tin nhắn trong hội thoại \\
\hline
\end{tabularx}
}


\end{document}